\section{Discussion}
The main point of the analysis is that three different notions of conditioning should not be
identified.  The Hermitian representative is governed by its eigenvalue ratio, the original
matrix is governed in Euclidean norm by its singular values, and the similarity transformation
is governed by the condition number of the chosen positive metric.  These quantities can track
one another, as in the two-level example, but they need not do so.

The three-level family isolates the failure sharply.  The defective limit occurs at the
nonzero eigenvalue $1$, so eigenvector coalescence forces every compatible positive metric to
become singular without driving the matrix toward noninvertibility.  The exact leading law
$\kappa_*(A_\varepsilon)\sim8/\varepsilon^2$ therefore measures a geometric obstruction to
constructing a well-conditioned physical inner product, not an instability of the Euclidean
inverse.  The normalized-solution estimate reinforces this distinction: small Euclidean
perturbations produce uniformly controlled solution-state changes despite the diverging metric
condition number.

The constant $8$ is also informative.  A bound based only on the overlap of the coalescing
pair captures the quadratic divergence but not the optimal leading constant.  The proof of the
sharp lower bound uses the simultaneous metric orthogonality of all three right eigenvectors,
and the third eigendirection remains relevant even though its eigenvalue stays separated.
This is consistent with the broader lesson from exceptional-point geometry that additional
states can influence eigenvector sensitivity beyond a two-state reduction, while our optimized
metric condition number is a distinct quantity.

The growing bidiagonal family shows that the phenomenon is not an artifact of fixed dimension.
For every $N\ge3$, the inverse norm remains bounded while the same $8/\varepsilon^2$ lower bound
is inherited from the first three eigendirections.  The Hermitian-dilation proposition then
serves as a useful diagnostic: its condition number is the ordinary singular-value condition
number, which remains bounded.  Thus neither the metric divergence nor the exceptional-point
limit alone yields a lower bound on quantum linear-system complexity.  Any such claim would
require a specified oracle, state-preparation model, precision criterion and readout task.

Several limitations should be kept explicit.  First, the exact finite-$\varepsilon$ minimizer
of $\kappa_2(\eta)$ for the three-level family is not characterized here; what is exact is the
leading asymptotic constant.  Second, the examples are structured finite-dimensional families,
not a classification of all pseudo-Hermitian exceptional points.  Third, perturbations are
measured in the reference Euclidean norm and need not preserve quasi-Hermiticity.  Finally,
metric singularity may have direct physical significance for observables or parameter-estimation
protocols even when Euclidean inversion remains stable; the present result concerns spectral
inversion specifically.
